\documentclass{article}

\usepackage{iftex}
\ifPDFTeX
  \usepackage[T1]{fontenc}
  \usepackage[utf8]{inputenc}
\fi
\usepackage[UTF8]{ctex}
\usepackage[backend=biber,style=gb7714-2015,sorting=gb7714-2015,gbalign=left,sorting=gb7714-2015]{biblatex}
\usepackage{amsmath, amsthm, amssymb, amsfonts}
\usepackage{thmtools}
\usepackage{graphicx}
\usepackage{setspace}
\usepackage{geometry}
\usepackage{float}
\usepackage{hyperref}
\usepackage[english]{babel}
\usepackage{framed}
\usepackage[dvipsnames]{xcolor}
\usepackage{tcolorbox}
\usepackage{tikz-cd}
\usepackage{tikz} 

\tcbuselibrary{breakable}

\addbibresource{references.bib}

\colorlet{LightGray}{White!90!Periwinkle}
\colorlet{LightOrange}{Orange!15}
\colorlet{LightGreen}{Green!15}

\newcommand{\HRule}[1]{\rule{\linewidth}{#1}}

\declaretheoremstyle[name=Theorem,]{thmsty}
\declaretheorem[style=thmsty,numberwithin=section]{theorem}
\tcolorboxenvironment{theorem}{colback=White}

\declaretheoremstyle[name=Definition,]{prosty}
\declaretheorem[style=prosty,numberlike=theorem]{definition}
\tcolorboxenvironment{definition}{colback=White}

\declaretheoremstyle[name=Lemma,]{prcpsty}
\declaretheorem[style=prcpsty,numberlike=theorem]{lemma}
\tcolorboxenvironment{lemma}{colback=White}

\newtheorem{proposition}[theorem]{Proposition}
\newtheorem{corollary}[theorem]{Corollary}
\newtheorem{remark}[theorem]{Remark}
\newtheorem{example}[theorem]{Example}
\newtheorem{exercise}[theorem]{Exercise}


\def\E{\mathop{\mathcal E}}
\def\D{\mathop{\mathcal D}}
\def\lhdeq{\mathop{\trianglelefteq}}
\def\<{\mathop{\langle}}
\def\>{\mathop{\rangle}}
\def\l{\mathop{\ell}}
\newcommand{\g}{\mathfrak{g}}  
\newcommand{\h}{\mathfrak{h}} 
\newcommand{\n}{\mathfrak{n}}  
\newcommand{\U}{\mathcal{U}}  
\newcommand{\C}{\mathbb{C}} 
\newcommand{\borel}{\mathfrak{b}}
\newcommand{\Ind}{\mathrm{Ind}} 
\newcommand{\Hom}{\mathrm{Hom}}  
\newcommand{\Ext}{\mathrm{Ext}}

\newcommand{\slC}{\mathfrak{sl}(2, \mathbb{C})}
\newcommand{\vir}{\mathfrak{vir}}
\newcommand{\der}{\mathrm{d}}

\newcommand{\Z}{\mathbb{Z}}
\newcommand{\End}{\mathrm{End}}
\newcommand{\N}{\mathbb{N}}

% ------------------------------------------------------------------------------

\begin{document}

% ------------------------------------------------------------------------------
% Cover Page and ToC
% ----------------------------------------------------------------------
\section*{讨论班安排}

\begin{itemize}
\item 时间：每周一9-10节（19:00-20：:50）。
\item 地点：一教121。
\item 指导教师：VProf. Vyacheslav Futorny。
\item 主持人：刘耀华。
\item 讨论班形式：自行报名，填写群文件中“2026春李代数讨论班pre报名”在线表格，在报名时间进行pre。
\item 主题要求：本讨论班是Zelmanov教授“李理论”课程的workshop，根据李理论课程内容自行选择topic进行讨论，讨论内容可以是课程内容的延伸和推广，也可以是参考书内容的介绍和课后习题解答。参考书为：Introduction to Lie Algebras and Representation Theory, J.E. Humphreys, Springer. 
\item 流程：周一，周三（双）为李理论课程时间，需要做pre的同学在这1-2次课中选取topic，在下周一的9-10节讨论班上进行pre，时长不超过90分钟。并且在14周前把你的讲稿发给主持人。讨论班结束后，主持人将对pre内容进行总结，并在群文件中上传总结内容。每次内容的记录（pre或纯讲稿，见下文）会记录贡献者姓名和pre时间，作为讨论班的参与记录。并不定时向Futorny教授或Zelmanov教授汇报记录内容，以展示讨论班的进展情况。
\item 学分相关：本讨论班为后置名单课程，请需要学分的同学qq私信我你的姓名与学号方便统计，我会在第14周将审核合格的名单上报给数学系完成学分评定。
\item 如何审核合格？
     \begin{enumerate}
         \item 至少完成一次pre，或者从做pre的同学没有涉及的内容中选取topic，提交一次讲稿。
         \item 积极参与线上，线下讨论。 
     \end{enumerate}
\item 补充：我会尽量邀请Futorny教授或Zelmanov教授，或者其他教授、博后来参加讨论班，听取同学们的pre内容，并给予指导。邀请成功后我会提前在群里通知大家。
\end{itemize}




\end{document}
